\documentclass[a4paper,12pt]{jsarticle}
\usepackage[margin=25mm]{geometry}
\usepackage{amsmath,amssymb}
\usepackage{enumitem}

\begin{document}

\section*{第1問（数と式・二次関数）}
\begin{enumerate}[label=\arabic*.]
  \item
  \[
  P(x)=-2(x-6)^2+32
  \]
  より、最大となるのは
  \[
  x=6.
  \]
  \item 最大利益は
  \[
  32\text{万円}.
  \]
  \item
  \[
  -2x^2+24x-40\ge0
  \iff x^2-12x+20\le0
  \iff (x-2)(x-10)\le0.
  \]
  よって
  \[
  2\le x\le10.
  \]
  \item
  \[
  P(3)=-2\cdot9+24\cdot3-40=14,
  \quad
  P(7)=-2\cdot49+24\cdot7-40=30.
  \]
  したがって増加量は
  \[
  30-14=16\text{万円}.
  \]
\end{enumerate}

\section*{第2問（場合の数・確率）}
\begin{enumerate}[label=\arabic*.]
  \item
  \[
  \binom{9}{3}=84.
  \]
  \item 奇数5枚、偶数4枚。和が偶数になるのは（奇奇偶）または（偶偶偶）。
  \[
  \binom52\binom41+\binom43=40+4=44.
  \]
  よって
  \[
  \frac{44}{84}=\frac{11}{21}.
  \]
  \item 最大が7なら7を含み、残り2枚は1～6から選ぶ。
  \[
  \binom62=15.
  \]
  よって
  \[
  \frac{15}{84}=\frac{5}{28}.
  \]
  \item $1\le a<b<c\le9$ で $a+b>c$ を満たす組を数えると26通り。よって
  \[
  \frac{26}{84}=\frac{13}{42}.
  \]
\end{enumerate}

\section*{第3問（図形と計量）}
\begin{enumerate}[label=\arabic*.]
  \item 余弦定理より
  \[
  BC^2=9^2+7^2-2\cdot9\cdot7\cos60^\circ
  =81+49-63=67.
  \]
  よって
  \[
  BC=\sqrt{67}.
  \]
  \item
  \[
  S=\frac12\cdot9\cdot7\sin60^\circ
  =\frac{63\sqrt3}{4}.
  \]
  \item 正弦定理より
  \[
  \frac{BC}{\sin60^\circ}=2R
  \Rightarrow
  R=\frac{\sqrt{67}}{\sqrt3}=\frac{\sqrt{201}}{3}.
  \]
  \item 高さを $h$ とすると
  \[
  \frac12\cdot BC\cdot h=S
  \Rightarrow
  h=\frac{2S}{BC}=\frac{63\sqrt3}{2\sqrt{67}}.
  \]
\end{enumerate}

\section*{第4問（データの分析）}
\begin{enumerate}[label=\arabic*.]
  \item
  \[
  H=50+10\frac{72-62}{10}=60.
  \]
  \item
  \[
  50+10\frac{x-62}{10}\ge60
  \iff x\ge72.
  \]
  \item
  \[
  H(52)=50+\frac{10(52-62)}{10}=40,
  \quad
  H(77)=50+\frac{10(77-62)}{10}=65.
  \]
  差は
  \[
  65-40=25.
  \]
  \item
  \[
  50+10\frac{x-62}{10}=45
  \iff x-62=-5
  \iff x=57.
  \]
\end{enumerate}

\end{document}
